% =====================================================================
% DRAFT v0.1 — OCR errors in Bangla/English NCTB math textbooks
% Status: v0.2 — Day 1 + Day 5 (CER on 3 verified Bangla pages). Every "TODO" must be filled from real
% experiments (Days 2-6). Do NOT invent numbers or citations.
% Later: move this content into the official ACL template on Overleaf.
% =====================================================================
\documentclass[11pt]{article}
\usepackage[margin=1in]{geometry}
\usepackage{booktabs}
\usepackage{amssymb}
\usepackage{hyperref}
\usepackage{xcolor}
\newcommand{\todo}[1]{\textcolor{red}{[TODO: #1]}}

\title{Where OCR Breaks on Bangla Mathematics Textbooks:\\
An Error Analysis of Tesseract, LLM Post-Correction, and Vision LLMs\\
for Retrieval-Augmented Math Tutoring}
\author{Peal Hasan \\ University of the People \\ \todo{email}}
\date{}

\begin{document}
\maketitle

% ---------------------------------------------------------------------
\begin{abstract}
Retrieval-augmented generation (RAG) tutors depend on clean text extracted
from source textbooks, yet the national (NCTB) mathematics textbooks used by
millions of secondary students in Bangladesh are distributed as scanned PDFs
that mix Bangla script, Latin variables and mathematical notation.
We study how optical character recognition (OCR) fails on these books and how
much of the damage can be repaired.
On 45 pages of a Grade 9--10 chapter on sets and functions (22 Bangla, 23
English), an LLM post-correction pass made 1{,}024 edits to Tesseract output;
52.6\% of them targeted mathematical notation (symbols, exponents and set
brackets) rather than ordinary words, and Bangla--Latin digit confusion
appeared only in Bangla pages.
On three hand-verified Bangla pages, Tesseract had a character error rate
(CER) of 11.0\% and misread 99\% of mathematical symbols; text-only LLM
post-correction lowered CER only to 8.4\% (symbols: 67\% wrong), whereas a
vision LLM reading the page image reached 0.24\% CER with no symbol errors.
\todo{1 sentence: effect on RAG answer accuracy for 20 questions (Day 6).}
We release our annotated pages and code to support Bangla educational NLP.
\end{abstract}

% ---------------------------------------------------------------------
\section{Introduction}
\begin{itemize}
  \item Problem: AI math tutors for Bangla-medium students need the official
        NCTB textbooks as grounding; the books exist only as scanned images.
  \item Why it is hard: mixed scripts (Bangla words, Latin variables, both
        Bangla and Latin digits) plus math notation ($\in, \cup, \cap,
        \emptyset, \therefore$, superscripts, set braces).
  \item Our system context: MathAI Bangladesh, a RAG tutor (Section~\ref{sec:system}).
        Observation that motivated this paper: answers were often irrelevant
        or wrong, and the cause traced back to OCR text.
  \item Contributions:
    \begin{enumerate}
      \item A categorised analysis of 1{,}024 OCR corrections on 45 bilingual
            textbook pages.
      \item A small character-level ground truth \todo{N pages} and a CER
            comparison of three extraction methods.
      \item Evidence that LLM post-correction is a double-edged tool: it fixes
            notation but also deletes or rewrites correct text.
      \item \todo{RAG answer accuracy before/after correction.}
    \end{enumerate}
\end{itemize}

% ---------------------------------------------------------------------
\section{Related Work}
\todo{Keep short (3 paragraphs). Only cite papers you have actually opened.}
\begin{itemize}
  \item OCR engines: Tesseract \cite{smith2007tesseract}. \todo{Bangla OCR
        work; math OCR / formula recognition work.}
  \item LLM-based OCR post-correction. \todo{find 2--3 papers.}
  \item Retrieval-augmented generation \cite{lewis2020rag}; effect of noisy
        retrieval corpora. \todo{find 1--2 papers.}
  \item Bangla NLP resources and educational AI. \todo{find 1--2 papers.}
\end{itemize}

% ---------------------------------------------------------------------
\section{System and Data}\label{sec:system}
\paragraph{Corpus.} Five NCTB mathematics textbooks (Grades 6--10, General and
Higher Math), each in a Bangla and an English edition: 10 scanned PDFs,
61 chapters, about 2{,}562 pages.

\paragraph{Pipeline.} Tesseract OCR (\texttt{ben+eng}) $\rightarrow$
chunking (1{,}000 characters, 200 overlap, sentence-aware) $\rightarrow$
multilingual-e5-large embeddings $\rightarrow$ Qdrant vector store
(3{,}907 chunks) $\rightarrow$ retrieval $\rightarrow$ answer generation with a
multi-LLM cross-checker (two models plus a tie-breaker).

\paragraph{Study subset.} Grade 9--10 General Math, Chapter 2 (Sets and
Functions): Bangla pages 26--47 and English pages 27--49 (45 pages,
60{,}521 raw OCR characters).

% ---------------------------------------------------------------------
\section{Method}
\paragraph{M1: Tesseract (baseline).} Raw output, \texttt{ben+eng}.
\paragraph{M2: Tesseract + LLM post-correction.} Each page's OCR text is sent to
a Groq-hosted LLM with a short Bangla prompt asking it to fix only
mathematical notation (e.g.\ \texttt{x2}$\rightarrow$x$^2$,
\texttt{...}$\rightarrow\therefore$); the edit list (before/after spans) is
logged per page. The exact model version used for this run could not be
recovered from the project logs; \todo{state as such or confirm}.
\paragraph{M3: Vision LLM transcription.} The page image (300 DPI) is
transcribed directly by \texttt{gemini-3.6-flash} (October 2026,
temperature 0) with a prompt asking for an exact, uncorrected transcription
with Unicode math symbols.
\paragraph{Error categories.} Each logged edit is assigned automatically to one
of seven categories (math symbol, superscript/exponent, bracket/set notation,
Bangla--Latin digit swap, punctuation/spacing, deletion, other) by
rule-based heuristics (code released).
\paragraph{Ground truth and metric.} Three Bangla pages (PDF pages 30, 36,
44; printed pages 25, 31, 39) covering definitions and worked examples (30,
44) and an exercise set (36) were verified character by character against
the page image by the author, a mathematics teacher. To save effort the
verification started from the vision-LLM transcription and corrected it;
this may favour M3, so further pages are being verified starting from the
Tesseract output instead (Section~\ref{sec:limits}). CER is the Levenshtein
distance divided by reference length after NFC normalisation and removal of
all whitespace (a variant keeping single spaces gives the same ranking).
We also break errors down by character class.
\paragraph{Downstream RAG test.} \todo{20 questions from Chapter 2 with
reference answers; answers marked correct/incorrect by the author.}

% ---------------------------------------------------------------------
\section{Results}
\subsection{What the LLM corrector changed}

\begin{table}[h]
\centering
\begin{tabular}{lrrr}
\toprule
 & Bangla & English & Total \\
\midrule
Pages & 22 & 23 & 45 \\
Raw OCR characters & 29{,}615 & 30{,}906 & 60{,}521 \\
LLM edits & 557 & 467 & 1{,}024 \\
Edits per page & 25.3 & 20.3 & 22.8 \\
\bottomrule
\end{tabular}
\caption{LLM post-correction edits on Chapter 2.}
\label{tab:edits}
\end{table}

\begin{table}[h]
\centering
\begin{tabular}{lrrrr}
\toprule
Category & Bangla & English & Total & \% \\
\midrule
Math symbol & 172 & 182 & 354 & 34.6 \\
Other (word / mixed) & 194 & 94 & 288 & 28.1 \\
Punctuation / spacing & 47 & 80 & 127 & 12.4 \\
Superscript / exponent & 46 & 79 & 125 & 12.2 \\
Bracket / set notation & 43 & 16 & 59 & 5.8 \\
Bangla--Latin digit swap & 43 & 0 & 43 & 4.2 \\
Deletion & 12 & 16 & 28 & 2.7 \\
\midrule
Total & 557 & 467 & 1{,}024 & 100 \\
\bottomrule
\end{tabular}
\caption{Edits by category (automatic, heuristic labels).}
\label{tab:categories}
\end{table}

Observations:
\begin{itemize}
  \item Mathematical notation (symbols, exponents, set brackets) accounts for
        52.6\% of edits: the dominant OCR failure is notation, not Bangla words.
  \item Bangla--Latin digit confusion (e.g.\ \texttt{x} read as Bangla five)
        occurs only on Bangla pages, a mixed-script failure specific to
        Bangla mathematics books.
  \item Some edits are harmful: correct text deleted, a written number word
        rewritten as a digit, plain digits turned into superscripts.
        On the verified pages, Bangla-letter errors \emph{increased} after
        correction (0.98\%$\rightarrow$1.47\%; Table~\ref{tab:class}): the
        corrector rewrites text it cannot see.
\end{itemize}

\subsection{Character error rate}
\begin{table}[h]
\centering
\begin{tabular}{lrrrr}
\toprule
Method & p.30 & p.36 & p.44 & All (micro) \\
\midrule
M1 Tesseract & 7.98 & 13.43 & 13.09 & 11.03 \\
M2 Tesseract + LLM correction & 6.53 & 10.48 & 9.10 & 8.40 \\
M3 Vision LLM & 0.07 & 0.11 & 0.61 & 0.24 \\
\bottomrule
\end{tabular}
\caption{CER (\%) on three verified Bangla pages (3{,}273 reference characters).
A second, independent Tesseract run gave 10.91\%.}
\label{tab:cer}
\end{table}

\begin{table}[h]
\centering
\begin{tabular}{lrrrr}
\toprule
Character class & $n$ & M1 & M2 & M3 \\
\midrule
Bangla letters / vowel signs & 1{,}836 & 0.98 & 1.47 & 0.11 \\
Bangla digits & 34 & 32.35 & 32.35 & 14.71 \\
Latin letters (variables) & 296 & 32.77 & 18.92 & 0.00 \\
ASCII digits & 291 & 0.69 & 2.41 & 0.00 \\
Math symbols & 99 & 98.99 & 66.67 & 0.00 \\
Brackets / punctuation & 717 & 14.92 & 7.53 & 0.00 \\
\bottomrule
\end{tabular}
\caption{Error rate (\%) by reference character class.}
\label{tab:class}
\end{table}

Tesseract reads Bangla prose well (under 1\% error on Bangla letters) but
fails almost completely on mathematical notation: 98 of 99 math symbols and a
third of the Latin variables were wrong. Text-only correction recovers some
notation but cannot restore what OCR has destroyed (e.g.\ a garbled
\texttt{(2,\%,2]} for $\{x,y,z\}$), and introduces new errors in Bangla text
and digits. The vision LLM's few errors were confined to Bangla digits in a
vertically printed margin year and one extra vowel sign in a page footer.

\subsection{Effect on RAG answers}
\todo{Accuracy on 20 questions with current (uncorrected) index; if time,
after re-indexing corrected text (Day 6).}

% ---------------------------------------------------------------------
\section{Discussion and Limitations}
\begin{itemize}
  \item Small scale: one chapter, three verified pages, one annotator.
        Ground truth started from the M3 output, which may favour M3.\label{sec:limits}
  \item Category labels are heuristic; a manual check of a sample is
        \todo{done / future work}.
  \item Models change quickly; results are tied to the model versions named.
  \item Future work: all 61 chapters, human-reviewed corrections, full RAG
        evaluation.
\end{itemize}

\section{Conclusion}
\todo{3--4 sentences after results are final.}

% ---------------------------------------------------------------------
\begin{thebibliography}{9}
\bibitem{smith2007tesseract}
R.~Smith. An overview of the Tesseract OCR engine. In \emph{Proc.\ ICDAR}, 2007.
\bibitem{lewis2020rag}
P.~Lewis et al. Retrieval-augmented generation for knowledge-intensive NLP
tasks. In \emph{Advances in NeurIPS}, 2020.
\end{thebibliography}

\end{document}
